\documentclass[11pt]{article}

\usepackage[T1]{fontenc}
\usepackage{amsmath,amssymb,amsthm}
\usepackage[margin=1in]{geometry}
\usepackage[hidelinks]{hyperref}
\setlength{\emergencystretch}{2em}

\newtheorem{theorem}{Theorem}
\newtheorem{lemma}[theorem]{Lemma}
\newtheorem{corollary}[theorem]{Corollary}
\newtheorem{conjecture}[theorem]{Conjecture}
\theoremstyle{remark}
\newtheorem{remark}[theorem]{Remark}

\newcommand{\R}{\mathbb R}

\title{Multirecord Spatial Import into Fast Type-II Carriers}
\author{John Robert Lawson and Codex}
\date{Research round ending 24 July 2026}

\begin{document}
\maketitle

\begin{abstract}
This round proves a signed spatial-transport theorem for the retained
energy-efficient branch of the exact \(q=4\) Type-II ledger.  A future
radius-\(R_j\) carrier ball contains fixed energy at its record time
\(t_j\), but the first-record weak-\(L^3\) ceiling makes the energy in the
same ball vanish at
\(t_{j-\lfloor j/4\rfloor}\).  The viscous cutoff correction across this
multirecord window is \(o(1)\).  The exact local kinetic-energy equality
therefore forces a fixed net weighted pressure-plus-kinetic energy import
through the future carrier annulus.  These events occur on a pairwise
disjoint subsequence of time intervals.  Their numerical sum diverges.  The
finite common \(L^1\) flux budget controls only the summably weighted totals
\(R_j|\mathsf X_j|\), so no contradiction or freshness statement is proved.
The Clay problem remains unsolved.
\end{abstract}

\section{Epistemic scope}

The statements below are analytic theorems under the displayed
smoothness, finite-energy, first-record, energy-efficient, retained-geometry,
and exact \(q=4\) hypotheses.  The preceding terminal-infrared theorem says
that the retained endpoint core is ultraviolet; only its local total-energy
floor is needed in the present proof.  The imported total-energy amount is
not identified with that Fourier detector.  The result is new within this
repository and awaits external mathematical review.  No claim of literature
novelty is made.

\section{First-record carrier data}

Let \(u\) solve
\[
 \partial_tu+u\cdot\nabla u=-\nabla p+\nu\Delta u,
 \qquad
 \nabla\cdot u=0
\]
smoothly on \(\R^3\times[0,T^*)\), with
\[
 \sup_{t<T^*}\frac12\|u(t)\|_2^2\le E_0.
\]
Let \(t_j\uparrow T^*\) be first record times for
\[
 m_j:=\|u(t_j)\|_{L^{3,\infty}},
\]
so that
\[
 \|u(t)\|_{L^{3,\infty}}\le m_j
 \qquad(0\le t\le t_j).
\]

At \(t_j\), choose a near-extremising amplitude layer
\[
 A_j:=\{a_j<|u(t_j)|\le2a_j\},
\]
and define
\[
 e_j:=\int_{A_j}|u(t_j)|^2\,dx,
 \qquad
 R_j:=|A_j|^{1/3}.
\]
In the energy-efficient branch,
\[
 0<c_0\le e_j\le2E_0,
 \qquad
 R_j\asymp\frac{e_j}{m_j^2}.
\]
In a partial or tight geometry cell, choose \(0<\eta<\theta\).
There are fixed \(A<\infty\), centres \(x_j\), and
\[
 \gamma:=\eta c_0>0
\]
such that
\begin{equation}
 \label{eq:final-floor}
 \boxed{
 \int_{B_{AR_j}(x_j)}|u(x,t_j)|^2\,dx\ge\gamma
 }
\end{equation}
eventually.

\section{Capacity of the future ball}

\begin{lemma}[Finite-volume weak-\(L^3\) capacity]
\label{lem:capacity}
For every finite-measure \(\Omega\subset\R^3\) and
\(f\in L^{3,\infty}(\R^3)\),
\[
 \boxed{
 \int_\Omega|f|^2\,dx
 \le
 C|\Omega|^{1/3}\|f\|_{L^{3,\infty}}^2.
 }
\]
\end{lemma}

\begin{proof}
Put \(K=\|f\|_{L^{3,\infty}}\).  The distribution function on \(\Omega\)
satisfies
\[
 |\Omega\cap\{|f|>\lambda\}|
 \le
 \min\{|\Omega|,K^3\lambda^{-3}\}.
\]
Split the layer-cake identity for \(\int_\Omega|f|^2\) at
\[
 \lambda_0:=K|\Omega|^{-1/3}.
\]
Then
\[
\begin{aligned}
 \int_\Omega|f|^2\,dx
 &=
 2\int_0^\infty
 \lambda|\Omega\cap\{|f|>\lambda\}|\,d\lambda\\
 &\le
 2|\Omega|\int_0^{\lambda_0}\lambda\,d\lambda
 +
 2K^3\int_{\lambda_0}^\infty\lambda^{-2}\,d\lambda\\
 &\le
 C K^2|\Omega|^{1/3}.
\end{aligned}
\]
\end{proof}

Choose earlier indices \(i_j<j\) such that
\begin{equation}
 \label{eq:amplitude-separation}
 \frac{m_{i_j}}{m_j}\longrightarrow0.
\end{equation}
At the earlier record, Lemma~\ref{lem:capacity}, the first-record property,
and \(R_j\lesssim E_0m_j^{-2}\) give
\begin{equation}
 \label{eq:initial-empty}
\begin{aligned}
 &\int_{B_{(A+1)R_j}(x_j)}
 |u(x,t_{i_j})|^2\,dx\\
 &\qquad\le
 C_AR_jm_{i_j}^2\\
 &\qquad\le
 C_AE_0\left(\frac{m_{i_j}}{m_j}\right)^2
 \longrightarrow0.
\end{aligned}
\end{equation}
This estimate is uniform in the future centres \(x_j\).

\section{The signed local energy identity}

Choose
\[
 0\le\chi\in C_c^\infty(\R^3),
 \qquad
 \chi=1\ \hbox{on }B_A,
 \qquad
 \operatorname{supp}\chi\subset B_{A+1},
\]
with \(\chi\) radial and nonincreasing in radius, and put
\[
 \chi_j(x):=
 \chi\left(\frac{x-x_j}{R_j}\right).
\]
Define
\[
 \mathcal E_j(t)
 :=
 \frac12\int_{\R^3}\chi_j|u(t)|^2\,dx.
\]

The smooth local energy equality is
\[
 \partial_t\frac{|u|^2}{2}
 +
 \nabla\cdot
 \left[
 \left(\frac{|u|^2}{2}+p\right)u
 \right]
 =
 \nu\Delta\frac{|u|^2}{2}
 -
 \nu|\nabla u|^2.
\]
Integrating it against \(\chi_j\) on \([t_{i_j},t_j]\) gives
\begin{equation}
 \label{eq:local-balance}
 \boxed{
 \mathcal E_j(t_j)-\mathcal E_j(t_{i_j})
 =
 \mathsf X_j-\mathsf D_j+\mathsf C_j,
 }
\end{equation}
where
\[
 \mathsf X_j
 :=
 \int_{t_{i_j}}^{t_j}\!\!\int_{\R^3}
 \left(\frac12|u|^2+p\right)
 u\cdot\nabla\chi_j\,dx\,dt,
\]
\[
 \mathsf D_j
 :=
 \nu\int_{t_{i_j}}^{t_j}\!\!\int_{\R^3}
 \chi_j|\nabla u|^2\,dx\,dt
 \ge0,
\]
and
\[
 \mathsf C_j
 :=
 \frac{\nu}{2}
 \int_{t_{i_j}}^{t_j}\!\!\int_{\R^3}
 |u|^2\Delta\chi_j\,dx\,dt.
\]
The pressure gauge is immaterial because
\(\int u\cdot\nabla\chi_j\,dx=0\).  With the sign in
\eqref{eq:local-balance}, positive \(\mathsf X_j\) is net energy import
through the annulus supporting \(\nabla\chi_j\).

Since
\[
 \|\Delta\chi_j\|_\infty
 =
 R_j^{-2}\|\Delta\chi\|_\infty,
\]
the global energy ceiling gives
\begin{equation}
 \label{eq:cutoff-error}
 \boxed{
 |\mathsf C_j|
 \le
 C_{\chi,E_0}\nu\frac{t_j-t_{i_j}}{R_j^2}.
 }
\end{equation}

\begin{theorem}[Multirecord retained-core import]
\label{thm:import}
Assume \eqref{eq:final-floor}, \eqref{eq:amplitude-separation}, and
\[
 \nu\frac{t_j-t_{i_j}}{R_j^2}\longrightarrow0.
\]
Then
\[
 \boxed{
 \liminf_{j\to\infty}\mathsf X_j
 \ge\frac{\gamma}{2}>0.
 }
\]
\end{theorem}

\begin{proof}
The final floor and the choice of cutoff give
\[
 \mathcal E_j(t_j)\ge\frac{\gamma}{2}.
\]
Equation~\eqref{eq:initial-empty} gives
\[
 \mathcal E_j(t_{i_j})\longrightarrow0,
\]
and \eqref{eq:cutoff-error} gives \(\mathsf C_j\to0\).
Rearranging \eqref{eq:local-balance},
\[
 \mathsf X_j
 =
 \mathcal E_j(t_j)-\mathcal E_j(t_{i_j})
 +
 \mathsf D_j-\mathsf C_j.
\]
The nonnegative dissipation \(\mathsf D_j\) can only increase the required
import.  Taking the lower limit proves the theorem.
\end{proof}

\begin{remark}
No absolute estimate on the pressure flux is used.  This is a lower bound
for its signed sum with kinetic transport.  The theorem does not assign
separate signs to the pressure and kinetic pieces.
\end{remark}

\section{The exact q-four schedule}

The representative ledger has, up to fixed comparison constants,
\[
 m_j\asymp2^{2j},
 \qquad
 R_j\asymp2^{-4j},
 \qquad
 t_{j+1}-t_j\asymp\frac{2^{-11j}}j.
\]
Set
\[
 N_j:=\left\lfloor\frac j4\right\rfloor,
 \qquad
 i_j:=j-N_j.
\]
Then
\[
 j-i_j=N_j\longrightarrow\infty
\]
and
\[
 \frac{m_{i_j}}{m_j}
 \asymp
 2^{-2N_j}
 \longrightarrow0.
\]
The geometric gap sum gives
\[
 t_j-t_{i_j}
 \le
 C\frac{2^{-11i_j}}{i_j}.
\]
Since \(R_j^2\asymp2^{-8j}\),
\[
\begin{aligned}
 \nu\frac{t_j-t_{i_j}}{R_j^2}
 &\le
 \frac{C\nu}{i_j}
 2^{-11(j-N_j)+8j}\\
 &=
 \frac{C\nu}{i_j}
 2^{-3j+11N_j}\\
 &\le
 \frac{C\nu}{j}
 2^{-j/4}
 \longrightarrow0.
\end{aligned}
\]
Here \(j\) remains the original q-four record-grid index.  A partial or
tight geometry subsequence is retained as an infinite subset of that grid
and is not renumbered.  Theorem~\ref{thm:import} therefore applies at every
sufficiently late retained index of any smooth trajectory realising the
exact survivor.

\begin{corollary}[Linear record depth]
\label{cor:depth}
For any fixed \(0<\alpha\le3/11\), the same conclusion holds with
\[
 i_j=j-\lfloor\alpha j\rfloor.
\]
Thus the future carrier ball is empty at an earlier record whose index
distance from \(j\) grows linearly.
\end{corollary}

\begin{proof}
The amplitude ratio is comparable to
\[
 2^{-2\lfloor\alpha j\rfloor}\longrightarrow0.
\]
The cutoff error is bounded by a fixed multiple of
\[
 j^{-1}2^{(-3+11\alpha)j},
\]
which tends to zero when \(\alpha<3/11\).  At \(\alpha=3/11\), the
exponential factor is bounded and the remaining \(1/i_j\) factor tends to
zero.
\end{proof}

\section{Disjoint event extraction}

For \(i_j=j-\lfloor j/4\rfloor\), choose \(j_n\) recursively so that
\[
 i_{j_{n+1}}>j_n.
\]
Then the intervals
\[
 [t_{i_{j_n}},t_{j_n}]
\]
are pairwise disjoint.  Theorem~\ref{thm:import} gives, eventually,
\[
 \boxed{
 \mathsf X_{j_n}\ge\frac{\gamma}{4}.
 }
\]
Hence
\[
 \sum_n\mathsf X_{j_n}=+\infty.
\]

\begin{remark}[What nonsummability means here]
This is a nonsummable ledger of signed event totals on one trajectory, not
a contradiction with a known finite Navier--Stokes action.  Each event uses
a different moving, shrinking cutoff.  A single packet can therefore be
counted each time it enters the next nested carrier ball.
\end{remark}

\section{The scale-weighted common flux budget}

Choose the whole-space pressure representative
\[
 p=\mathcal R_a\mathcal R_b(u_au_b).
\]
Sobolev interpolation gives
\[
 \|u(t)\|_3^3
 \le
 C\|u(t)\|_2^{3/2}\|\nabla u(t)\|_2^{3/2}.
\]
The energy equality and H{\"o}lder in time imply
\[
 \int_0^{T^*}\|u(t)\|_3^3\,dt<\infty.
\]
Moreover,
\[
 \|p(t)\|_{3/2}\le C\|u(t)\|_3^2.
\]
Hence
\[
 \boxed{
 \mathcal F_*:=
 \int_0^{T^*}\!\!\int_{\R^3}
 \left(|u|^3+|p||u|\right)\,dx\,dt
 <\infty.
 }
\]

Since
\[
 \|\nabla\chi_j\|_\infty\le\frac{C_\chi}{R_j},
\]
the pairwise disjoint event intervals give
\[
 \boxed{
 \sum_nR_{j_n}|\mathsf X_{j_n}|
 \le
 C_\chi\mathcal F_*.
 }
\]
This is compatible with fixed \(\mathsf X_{j_n}\), because the exact carrier
radii have a finite geometric sum.  The direct energy-class transport
estimate therefore introduces one carrier-radius weight.

\section{Robust findings and open obligations}

\subsection*{Findings, subject to review}

\begin{enumerate}
\item A future carrier ball has vanishing energy at a record a linear number
of \(q=4\) generations earlier.
\item The retained endpoint has a fixed physical energy floor.
\item Viscous diffusion through the cutoff is negligible across the
multirecord window.
\item The exact local energy equality forces fixed positive signed
pressure-plus-kinetic import into the future carrier ball.
\item The preceding terminal-infrared theorem locates the endpoint floor in
an ultraviolet core.
\item Infinitely many fixed import events occur on pairwise disjoint time
intervals of one trajectory.
\item The finite common \(L^1\) flux budget controls only
\(R_j|\mathsf X_j|\), whose forced lower weights are summable.
\end{enumerate}

\subsection*{Things still to prove}

\begin{enumerate}
\item Assign finite crossing multiplicity to energy-transport paths, or
control the centre and nesting geometry of the moving cutoffs.
\item Remove the summable carrier-radius weight from the finite global
pressure-plus-kinetic transport action.
\item Prevent one energy packet from paying every nested import event.
\item Prove post-entry residence or another non-reuse cost.
\item Exclude the diffuse and divergent-normalised-energy branches.
\end{enumerate}

\begin{conjecture}[Finite-multiplicity import]
Along every retained first-record Type-II sequence, either its positive
imports admit an assignment to energy-transport paths with uniformly finite
crossing multiplicity, or repeated nesting forces a nonsummable viscous,
pressure, or Lagrangian deformation action.
\end{conjecture}

No part of the conjecture is proved here.  This round proves neither
regularity nor breakdown and establishes no Clay alternative A--D.

\end{document}
